\properlane{patience-and-conversion}{Charity preserves people who may yet become wheat}
\subsection{A remedy that could injure the crop}
The servants propose to uproot what the enemy has sown. Their concern
for the field is evident in their first question: surely the owner had
sown good seed. Chrysostom treats their readiness as affection for the
crop, although their remedy lacks discretion. They ask permission; the
owner's answer corrects the remedy without denying their concern.
A person can desire the preservation of the good and still propose an
action that harms it.\footnote{Chrysostom, \work{Homilies on Matthew},
XLVI.1, on Matt 13:24--30; NPNF I.10.}

The master's reason is expressed in the appointed words about rooting
up wheat together with tares. The danger concerns the means of separation.
The crop has grown together in one field, and removal undertaken now
cannot be assumed harmless to the wheat. Jesus's explanation gives this
danger a human subject: the seed signifies people. Chrysostom consequently
asks what happens to people when the servants' proposed removal becomes
violent action against error.

He first identifies the tares with heretical societies. Error appears
alongside truth and resembles it, as tares resemble wheat. He compares
false prophets after the prophets and false apostles after the apostles.
This resemblance explains deception, while the order of sowing preserves
the priority of the good. The enemy does not possess the field or supply
an alternative good creation. He corrupts what has been planted.

Chrysostom also takes sleep as a warning to those entrusted with keeping
the field. The negligence is moral, so wakefulness is possible. Leaders
and hearers alike need it. This creates a difficulty for anyone who would
read the owner's refusal as permission to ignore teaching: the very
homily that rejects killing calls the hearers to vigilance against error.
Its two directions concern different actions.

His application of uprooting has two explanations. An armed campaign
against heretics can bring war and kill saints along with the guilty.
The good suffers from the chosen remedy. He then considers the guilty
persons themselves: some tares can change and become wheat. Killing
them would destroy people who still have room for repentance. The harm
to wheat thus includes both persons already good and persons who may
become good. The latter possibility explains why a judgement about
someone's present error cannot settle his whole future.

Chrysostom is specific about the action he forbids here. He rejects
killing while allowing the checking of heretical speech and the breaking
up of assemblies. That distinction belongs to his actual argument. His
explanation can support restraint against destructive violence and hope
for conversion; it cannot be cited as unrestricted modern toleration.
Nor does he erase final punishment. If the tares remain unchanged at
the end, the owner orders the burning. Waiting preserves an opportunity
that the harvest will close.

\subsection{Augustine addresses the hearers themselves}
Augustine's sermon brings the warning into the Christian assembly.
He addresses evil Christians as tares and commands them to change before
harvest. He also addresses good Christians who suffer the company of
the wicked. Their present mingling in the field will end at the barn.
These two addresses prevent a hearer from treating the parable solely
as a classification of outsiders.\footnote{Augustine, Sermon XXIII
[Benedictine LXXIII], §§1--4, NPNF I.6.}

Augustine explains why a plant image can include human conversion.
Parables use several visible things to signify a person or a spiritual
reality. The bad soil in the previous parable and the tares in this one
can both signify wicked Christians. The image does not require a human
being to be as unchangeable as a botanical species. Augustine therefore
exhorts the hard ground to be ploughed and the tares to become wheat.
His appeal is addressed to persons capable of hearing and repentance.

There is also a difference between the servants and the final reapers.
The angels make no mistakes in the separation, whereas human beings
can mistake wheat for tares. Zeal cannot confer that missing certainty.
An eager person can make a waste of the very harvest he wants to save.
Augustine consequently tells the good to heed the warning that one who
thinks himself standing should take care lest he fall. The question is
both whether they identify others correctly and whether they will
persevere themselves.

He includes clergy and laity in this examination. Wheat and tares exist
among those in the high seats as well as among those below them. Office
does not provide an exemption from conversion, and membership among the
laity does not imply wickedness. The hearer cannot identify a permanently
pure group by a visible place within the church. Augustine's final appeal
asks the bad to imitate the good and all to attain God through mercy.

Chrysostom's immediate subject is heresy; Augustine's is wicked Christian
conduct more generally. Their common argument concerns the injury caused
by premature destructive action and the possibility of conversion.
Augustine adds a developed account of mistaken human identification.
These explanations make the master's reason more precise than a general
instruction to remain calm. There are people whose life or repentance
the impatient remedy would foreclose.

\subsection{The community that bears and corrects}
Colossians names what happens between a grievance and a possible response.
A person has a complaint against another; he is told to bear with and
forgive him as the Lord forgave. Chrysostom's exposition describes a
contest within the injured person's thoughts. One thought proposes
revenge; the other proposes endurance. Peace acts as umpire and awards
the prize to endurance because God has made peace with those who now
need to imitate his mercy.\footnote{Chrysostom, \work{Homilies on
Colossians}, VIII, on Col 3:13--15.}

The umpire image assigns judgement to peace rather than anger. Anger
may still make its case, but it does not determine the response. The
reason is the Lord's freely given reconciliation, which the believer
has already received. Chrysostom connects this with Paul's command to
be thankful: the forgiven servant's treatment of another servant reveals
whether he acknowledges what his own Lord has done for him.

Placed beside the Gospel, this account describes a practice available
to those who have been told to let the crop grow. The servants' restraint
can be inhabited by active charity. They can forgive an actual wrong,
teach truth and seek repentance. Paul's adjoining commands to teach
and admonish forbid reducing the interval to mere avoidance of contact.
There is a neighbour to be instructed and loved while his future remains
open. The act of instruction must itself answer to the mercy Paul
commands.

The Collect names these people God's family and asks its continual
preservation. Its reliance on heavenly grace explains why correction
cannot begin from the assumption that the corrector has secured himself.
The Secret is still more personal in its acknowledgement: the offerers
have offences to be pardoned and hearts that waver. Read in this setting,
their prayer concerns the very people who may have judged a neighbour's
fault accurately. Accurate recognition of his fault does not remove
their own need for absolution and guidance.

The cry from Psalm 129, heard in both Alleluia and Offertory, similarly
places the petitioner among those needing to be heard. Augustine's
question about who could stand under a strict reckoning of sin allows
no class of human hearers to claim the final reaper's innocence.
The Communion's confidence likewise belongs to people commanded by
Christ to forgive while they pray. Bede's account of charitable faith
keeps that relation to an offender within the act of petition. The
assembly's prayer for mercy and its duty to show mercy concern the same
persons.

\subsection{Peace, thanksgiving and the end of patience}
The Introit's promise of peace addresses exiles who must wait and seek
the good of the city where they live. Jerome's explanation of their
false hopes makes the interval concrete: demanding an immediate return
does not fulfil God's promise. It is possible to obey during the delay.
In this interpretation, the comparison with the Gospel concerns that
obedience while God's final action is awaited. The exiles themselves
are the historical recipients of Jeremias's letter; their neighbours
are not reclassified as the Gospel's weeds.

The Gradual ascribes rescue to God. In Augustine's exposition of
Psalm 43, persecuted Christians can praise a deliverance still awaited
with prophetic certainty. Their enemies' present power does not make
God changeable. In relation to the parable, this gives a reason to
await the owner's judgement rather than demand that private vengeance
make it visible now. The lament's suffering remains real; so does the
promise of divine rescue.

Thanksgiving in the Preface worships the Father, Son and Holy Spirit.
Paul's one body gives thanks through Christ, and the Preface expressly
praises the triune God. A community learning concord therefore has a
common act beyond the settlement of its quarrels. It acknowledges God
as the source of the mercy from which its members live.

After Communion the people ask for salvation's effect, having received
its pledge. That request can be heard beside the harvest: a gift has
begun whose fulfilment still lies ahead. Sacramental reception is a
present gift; the final verdict remains Christ's. Augustine's warning
to the supposedly secure hearer and the Postcommunion's continuing
petition both concern people who must persevere.

The barn and fire make the end morally serious. The good are gathered,
and persistent evil is judged. Patience does not require indifference
to that difference. Chrysostom and Augustine explain why hope of
conversion and care for the innocent must affect action before the end.
A person who can still repent is among those whom destructive haste
can injure.

\subsection{Four senses}
\begin{fourSenses}
\item[Literal.] The owner postpones uprooting because it can injure wheat,
then commands a separation at harvest. Paul requires forgiven people to
bear with, forgive, teach and admonish one another. The prayers ask grace,
pardon and protection for that same community.
\item[Allegorical.] Christ sows the children of the kingdom and forms one
body through his mercy. Chrysostom's possible future wheat and Augustine's
changing wicked Christians express conversion within Christ's present
work. The Church's holiness comes from him amid the present mixture.
\item[Moral.] The injured person lets Christ's peace decide between revenge
and endurance. Correction seeks the neighbour's conversion, recognises
the danger of destroying the innocent or reformable, and begins with
one's own need for forgiveness and guidance.
\item[Anagogical.] Christ's angels finally separate with the certainty
human servants lack. The barn promises gathering, while the fire warns
against persisting in evil. The saving effect requested after the
sacramental pledge concerns the community's arrival at that final end.
\end{fourSenses}
